\documentclass[11pt,a4paper]{scrbook}

\usepackage[english]{babel}
\usepackage{fontspec}
\usepackage{microtype}
\usepackage[a4paper,left=22mm,right=55mm,top=25mm,bottom=27mm,marginparwidth=40mm,marginparsep=8mm]{geometry}
\usepackage{marginnote}
\usepackage{xcolor}
\usepackage{enumitem}
\usepackage{booktabs}
\usepackage[hidelinks]{hyperref}

\IfFontExistsTF{Libertinus Serif}{%
  \setmainfont{Libertinus Serif}
  \setsansfont{Libertinus Sans}[Scale=0.94]
}{%
  \setmainfont{Latin Modern Roman}
  \setsansfont{Latin Modern Sans}
}

\definecolor{Accent}{HTML}{35586C}
\definecolor{Muted}{HTML}{666666}

\setkomafont{chapter}{\sffamily\bfseries\Huge}
\setkomafont{section}{\sffamily\bfseries\Large}
\setkomafont{subsection}{\sffamily\bfseries\large}

\newcommand{\marginlabel}[1]{\marginnote{\footnotesize\sffamily\bfseries\color{Accent}#1}}
\newenvironment{definition}[1]
  {\par\medskip\noindent{\sffamily\bfseries\color{Accent}Definition — #1}\par\smallskip}
  {\par\medskip}
\newenvironment{example}
  {\par\medskip\noindent{\sffamily\bfseries Example}\par\smallskip}
  {\par\medskip}
\newenvironment{exercise}[1]
  {\par\medskip\noindent{\sffamily\bfseries Exercise #1}\par\smallskip}
  {\par\medskip}

\begin{document}

\frontmatter
\title{Kao Textbook Reference}
\subtitle{A compact 1.5-column educational architecture}
\author{Author}
\maketitle
\tableofcontents

\mainmatter
\chapter{A learning chapter}
\marginlabel{Chapter map}
The main column carries the teaching sequence. The outer margin carries compact contextual
information, references, definitions, navigational cues, or supporting examples.

\section{Concept}
\begin{definition}{Active information margin}
A margin is active when it reduces interruption by holding useful secondary information
beside the passage it supports.
\end{definition}

\marginlabel{Remember}
Use recurring semantic components consistently. A textbook should feel like one system,
not a sequence of individually styled pages.

\section{Worked example}
\begin{example}
State the principle, show one representative case, then transition into learner practice.
\end{example}

\section{Practice}
\begin{exercise}{1}
Apply the concept to a new example. Leave enough space for the expected response modality.
\end{exercise}

\end{document}
